% TOP_PROBLEM: {"id": 1141, "title": "The Weinstein Conjecture", "category_id": 11, "category": {"id": 11, "name": "dynamical_systems", "display_name": "Dynamical Systems", "description": "Problems about long-term behavior of deterministic systems, Hamiltonian dynamics, and periodic orbits.", "slug": "dynamical-systems", "order_index": 11, "created_at": "2026-07-31T15:26:25.671Z"}, "status": "open"}
% FIRST_POSTED: 2026-09-25
% ENTRY_KIND: statement_only
% !TeX program = lualatex
% Definition and sources only; this file is not an LLM solution attempt.
\documentclass[11pt]{article}
\usepackage[margin=25mm]{geometry}
\usepackage{fontspec}
\setmainfont{Latin Modern Roman}
\usepackage{amsmath,amssymb,mathtools}
\usepackage{unicode-math}
\setmathfont{Latin Modern Math}
\usepackage{xurl,hyperref}
\hypersetup{colorlinks=true,urlcolor=blue}
\setcounter{secnumdepth}{0}
\setlength{\parindent}{0pt}
\setlength{\parskip}{5pt}
\setlength{\emergencystretch}{3em}
\begin{document}
\section*{\texorpdfstring{The Weinstein Conjecture}{The Weinstein Conjecture}}\label{1141}
\subsection{Definitions and mathematical statement}
On a closed smooth $(2n+1)$-manifold $M$, a contact form is a one-form $\alpha$ with $\alpha\wedge({\,\mathrm{d}}\alpha)^n$ nowhere zero. Its Reeb vector field $R_\alpha$ is uniquely specified by
\[
 \alpha(R_\alpha)=1,\qquad \iota_{R_\alpha}{\,\mathrm{d}}\alpha=0.
\]
A periodic orbit is a smooth map $\gamma:{\mathbb{R}}/T{\mathbb{Z}}\to M$, $T>0$, satisfying $\dot\gamma=R_\alpha(\gamma)$. The conjecture asserts the existence of such an orbit for every closed contact manifold and every contact form in all dimensions. Here closed means compact without boundary. No nondegeneracy, contractibility, or uniform bound on the period is imposed.
\par\smallskip\textit{Formulation and concept references:} \href{https://arxiv.org/pdf/2005.09568}{[S1]}.

\subsection{Short English statement}
Must every Reeb flow on a compact contact manifold without boundary have at least one closed trajectory?
\subsection{Sources}
\begin{itemize}
\item[S1]Eva Miranda and Cedric Oms, The singular Weinstein conjecture, Advances in Mathematics 389 (2021), Section 5.1, Conjecture 5.1.
\url{https://arxiv.org/pdf/2005.09568}.
\textit{Formulation source.}
\end{itemize}
\end{document}
